% generated from papers/XXI-GL2/paper_XXI_GL2.tex -- edit the paper, then rerun assemble.py
\chapter[The gap is exactly one]{Localized $GL_2$ Bost--Connes systems:\\ the equilibrium range and the automorphic obstruction range\\ are disjoint, by exactly one}\label{ch:XXI}
\chapterorigin{This chapter is Paper XXI of the series (\texttt{papers/XXI-GL2/}).}
\begin{summary}
We examine the proposal to characterize Langlands reciprocity for weight-two forms through
the $\mathrm{KMS}$ structure of a localized $GL_2$ Bost--Connes system, and find a
quantitative obstruction that closes it.

The Hilbert--Schmidt identity is correct: with $U_p=p^{-1/2}\rho_f(\Frob_p)$ unitary by
Deligne, $\|I_2-U_p\|_\HS^2=4(1-\cos\theta_p)=8\sin^2(\theta_p/2)$, and
$\Xi_1(\pi_f,P_\Q)=\infty$ for every cuspidal eigenform, by Jacquet--Shalika. But the
$L$-function formula proposed for it is wrong: $\Xi_\beta$ is \emph{linear} in $\Tr U_p$, so
only the standard $L$-function appears,
\[
\Xi_\beta(\pi_f,S)=4\log\zeta_S(\beta)-2\,\Re\log L_S(\beta,\pi_f)+O(1),
\]
matching Theorem 4.1 of Chapter~\ref{ch:VI} at $n=2$; the term
$-2\Re\log L_S(\beta,\Sym^2\pi_f)$ is
spurious. Rankin--Selberg governs the second moment $\sum_p|\Tr U_p|^2p^{-s}$, which is
Sato--Tate, not $\Xi_\beta$.

The decisive obstruction is a mismatch of temperature ranges, and it is exact. For any
$S$, the automorphic obstruction can be nontrivial only where $\sum_{p\in S}p^{-\beta}$
diverges, since $\Xi_\beta\le8\sum_{p\in S}p^{-\beta}$; write $\beta_c(S)\le1$ for that
abscissa. The $GL_2$ partition function has local factor
$(1-p^{-\beta})^{-1}(1-p^{1-\beta})^{-1}$ --- the Solomon zeta of $M_2(\Z_p)$ --- so it
converges exactly for $\beta>1+\beta_c(S)$. Hence
\[
\text{obstruction nontrivial}\iff\beta\le\beta_c(S),
\qquad
\text{Gibbs states exist}\iff\beta>1+\beta_c(S),
\]
and the two regions are disjoint with a gap of width exactly $1$, for every $S$ and every
automorphic $\pi$. For $S=P_\Q$: obstruction on $(0,1]$, equilibrium on $(2,\infty)$. There is
no temperature at which the $GL_2$ system has equilibrium states and the automorphic
obstruction can break symmetry.

The shift by one is not accidental. The number of sublattices of index $p$ in $\Z_p^2$ is
$p+1$ while the norm is $p$, and by Theorem 1.2 of Chapter~\ref{ch:XII} this excess is universal:
$d_\pp=(q^d-1)/(q-1)>q$. The same off-by-one of $\mathbb P^{d-1}$ that empties the boundary
quotient there displaces the equilibrium range here.

Finally, the proposed algebra is not available as stated: $M_2(\Z)^+$ has infinite unit group
$SL_2(\Z)$, so by the unit-group exclusion of Chapter~\ref{ch:X} (its Remark 1.5) its partition function
diverges identically and one
must pass to the Hecke quotient, which is not a semigroup crossed product.
\end{summary}


\section{The algebra}

The framework is that of [Ch.\begin{proposition}[The semigroup crossed product is unavailable]\label{XXI:prop:units}
Let $P=M_2(\Z)\cap GL_2^+(\Q)$. Then $P^\times=SL_2(\Z)$ --- an element of $P$ with inverse
in $P$ has positive integral determinant on both sides, hence determinant $1$ --- is
infinite and $|\det|=1$ on it, so
$\sum_{g\in P}|\det g|^{-\beta}\ge|P^\times|=\infty$ for every $\beta$. Hence
$C(Y)\rtimes P$ carries no Gibbs state at any temperature, and the construction must be
replaced by the Hecke-algebra quotient of Connes--Marcolli \cite{ConnesMarcolliBook}, in which one sums over
$GL_2(\Z)$-orbits.
\end{proposition}

\begin{proof}
This is the unit-group exclusion recorded in Chapter~\ref{ch:X} (its Remark 1.5): a multiplicative norm
is $1$ on units, so an infinite unit group makes
$\zeta_P$ diverge identically.
\end{proof}

\begin{proposition}[The Solomon zeta]\label{XXI:prop:solomon}
Summing over $GL_2(\Z)$-orbits, the local factor at $p$ is the Solomon zeta of $M_2(\Z_p)$
\cite{Solomon},
\[
\zeta_{GL_2,p}(\beta)=\bigl(1-p^{-\beta}\bigr)^{-1}\bigl(1-p^{1-\beta}\bigr)^{-1},
\]
so for a set $S$ of primes
\begin{equation}\label{XXI:eq:zetaGL2}
\zeta_{GL_2,S}(\beta)=\prod_{p\in S}\bigl(1-p^{-\beta}\bigr)^{-1}\bigl(1-p^{1-\beta}\bigr)^{-1},
\end{equation}
which converges if and only if $\sum_{p\in S}p^{1-\beta}<\infty$.
\end{proposition}

\begin{proof}
The Gibbs weight of the Hecke quotient sums one term per \emph{left} coset $GL_2(\Z)g$,
equivalently per sublattice $L\subseteq\Z_p^2$, weighted by its index \cite{ConnesMarcolliBook} --- not one
term per double coset. The number of sublattices of $\Z_p^2$ of index $p^n$ is
$\sigma(p^n)=1+p+\cdots+p^n$, and
\[
\sum_{n\ge0}\sigma(p^n)\,p^{-\beta n}
=\sum_{n\ge0}\frac{p^{n+1}-1}{p-1}\,p^{-\beta n}
=\bigl(1-p^{-\beta}\bigr)^{-1}\bigl(1-p^{1-\beta}\bigr)^{-1}.
\]
Convergence of the Euler product
is convergence of $\sum_p(p^{-\beta}+p^{1-\beta})$, dominated by the second term.
\end{proof}

\section{The obstruction series and its $L$-function}

\begin{lemma}[Hilbert--Schmidt identity]\label{XXI:lem:HS}
Let $f\in S_2(\Gamma_0(N))$ be a normalized cuspidal Hecke eigenform, $p\nmid N$, and
$U_p=p^{-1/2}\rho_f(\Frob_p)$, unitary with eigenvalues $e^{\pm i\theta_p}$ by Deligne's
theorem \cite{Deligne}. Then
\[
\bigl\|I_2-U_p\bigr\|_\HS^2=2\bigl(2-\Re\Tr U_p\bigr)=4\bigl(1-\cos\theta_p\bigr)=8\sin^2(\theta_p/2),
\]
so that
\begin{equation}\label{XXI:eq:XiGL2}
\Xi_\beta(\pi_f,S)=8\sum_{p\in S}p^{-\beta}\sin^2\!\bigl(\theta_p/2\bigr).
\end{equation}
\end{lemma}

\begin{proof}
Lemma 2.2 of Chapter~\ref{ch:VI} (its Trace form) with $n=2$, together with $\Tr U_p=a_p/\sqrt p=2\cos\theta_p$.
\end{proof}

\begin{theorem}[The $L$-function dictionary, corrected]\label{XXI:thm:L}
For $\beta>1/2$,
\begin{equation}\label{XXI:eq:Lcorrect}
\Xi_\beta(\pi_f,S)=4\log\zeta_S(\beta)-2\,\Re\log L_S(\beta,\pi_f)+O(1),
\end{equation}
with $L_S(s,\pi_f)=\prod_{p\in S}\bigl((1-e^{i\theta_p}p^{-s})(1-e^{-i\theta_p}p^{-s})\bigr)^{-1}$
the analytically normalized standard $L$-function. No symmetric square factor occurs.
\end{theorem}

\begin{proof}
By \eqref{XXI:eq:XiGL2}, $\Xi_\beta=4\sum_pp^{-\beta}-2\sum_pp^{-\beta}\Tr U_p$. The first sum is
$\log\zeta_S(\beta)+O(1)$; the second is $\Re\log L_S(\beta,\pi_f)+O(1)$, the prime-power
terms converging for $\beta>1/2$. This is Theorem 4.1 of Chapter~\ref{ch:VI} at $n=2$, where the general
form is $2n\log\zeta_{K,S}(\beta)-2\Re\log L_S(\beta,\rho)+O(1)$.
\end{proof}

\begin{remark}[Why $\Sym^2$ cannot appear]\label{XXI:rem:sym2}
The proposed formula adds $-2\Re\log L_S(\beta,\Sym^2\pi_f)$. That term cannot be present:
$\|I-U\|_\HS^2=2(2-\Re\Tr U)$ is \emph{linear} in $\Tr U_p$, whereas $\Sym^2$ enters through
$\Tr(U_p^2)=\Tr(U_p)^2-2\det U_p$, a quadratic expression. Rankin--Selberg does govern the
second moment: $\log L_S(s,\pi\times\pi^\vee)=\sum_p|\Tr U_p|^2p^{-s}+O(1)$, whose pole at
$s=1$ is the input to Sato--Tate. That is a different series from $\Xi_\beta$.
\end{remark}

\begin{corollary}[Divergence at $\beta=1$]\label{XXI:cor:diverge}
For $S=P_\Q$ and any cuspidal eigenform $f$, $\Xi_1(\pi_f,P_\Q)=\infty$.
\end{corollary}

\begin{proof}
As $\beta\downarrow1$, $\log\zeta(\beta)\to\infty$ while $L(\beta,\pi_f)$ has neither zero nor
pole at $\beta=1$ --- Jacquet--Shalika non-vanishing on $\Re s=1$ \cite{JS}, and cuspidality excluding a
pole --- so $\Re\log L_S(\beta,\pi_f)=O(1)$. Apply \eqref{XXI:eq:Lcorrect}.
\end{proof}

\section{The gap is exactly one}

\begin{definition}\label{XXI:def:betac}
For $S$ a set of primes let $\beta_c(S)=\inf\{\beta:\sum_{p\in S}p^{-\beta}<\infty\}$, so
$0\le\beta_c(S)\le1$.
\end{definition}

\begin{theorem}[Disjoint ranges]\label{XXI:thm:gap}
For every set $S$ of primes and every automorphic representation $\pi$ of $GL_2$ unramified
outside $S$:
\begin{enumerate}
\item $\Xi_\beta(\pi,S)<\infty$ for every $\beta>\beta_c(S)$, so no representation is an
obstruction there; the obstruction can be nontrivial only for $\beta\le\beta_c(S)$;
\item $\zeta_{GL_2,S}(\beta)<\infty$ for $\beta>1+\beta_c(S)$ and $=\infty$ for
$\beta<1+\beta_c(S)$ --- at the boundary $\beta=1+\beta_c(S)$ itself, finite exactly when
$\sum_{p\in S}p^{-\beta_c(S)}$ converges --- so Gibbs states of
the $GL_2$ system exist only for $\beta\ge1+\beta_c(S)$, and only beyond it when the
critical series diverges.
\end{enumerate}
Consequently
\[
\bigl(0,\beta_c(S)\bigr]\ \cap\ \bigl[1+\beta_c(S),\infty\bigr)=\emptyset ,
\]
and the open interval $(\beta_c(S),1+\beta_c(S))$, of width exactly $1$, lies in neither
region. For $S=P_\Q$ the obstruction lives
on $(0,1]$ and the equilibrium states on $(2,\infty)$.
\end{theorem}

\begin{proof}
(1) By \eqref{XXI:eq:XiGL2}, $\Xi_\beta(\pi,S)\le8\sum_{p\in S}p^{-\beta}$, finite for
$\beta>\beta_c(S)$. (2) By Proposition~\ref{XXI:prop:solomon}, convergence is that of
$\sum_{p\in S}p^{1-\beta}=\sum_{p\in S}p^{-(\beta-1)}$, which for $\beta-1>\beta_c(S)$
converges, for $\beta-1<\beta_c(S)$ diverges, and at $\beta-1=\beta_c(S)$ is the critical
series itself.
\end{proof}

\begin{corollary}[The programme cannot close]\label{XXI:cor:cannot}
There is no temperature at which the localized $GL_2$ system possesses equilibrium states and
the automorphic obstruction $\Xi_\beta(\pi,S)$ is capable of producing symmetry breaking. In
particular the proposed uniqueness theorem on $(0,1]$ concerns a range in which the $GL_2$
system has no Gibbs states, and Corollary~\ref{XXI:cor:diverge}, though true, has no
$\mathrm{KMS}$ consequence for this system.
\end{corollary}

\begin{remark}[Numerical illustration]\label{XXI:rem:num}
For $S=P_\Q$: $\sum_{p\le10^6}p^{-\beta}$ equals $4.79$, $2.89$, $1.99$ at
$\beta=0.9,1.0,1.1$, the first two growing without bound and the third converging; while
$\zeta(\beta)\zeta(\beta-1)$ is divergent at $\beta\le2$ and equals $8.33$, $1.98$, $1.30$ at
$\beta=2.2,3,4$. The window $(1,2]$ has neither.
\end{remark}

\begin{remark}[The shift is the multiplicity--norm mismatch]\label{XXI:rem:mismatch}
The displacement by $1$ is the phenomenon isolated in Theorem 1.2 of Chapter~\ref{ch:XII}. There, for a
maximal order in a central simple algebra of degree $d$, the number of right ideals of reduced
norm $\pp$ is $d_\pp=(q^d-1)/(q-1)=1+q+\cdots+q^{d-1}$, strictly larger than the norm
$q=\Ns\pp$; for $d=2$ the excess is exactly the point at infinity of $\mathbb P^1$. Here the
same count --- $p+1$ sublattices of index $p$ in $\Z_p^2$ against norm $p$ --- produces the
factor $(1-p^{1-\beta})^{-1}$ in \eqref{XXI:eq:zetaGL2} and shifts the abscissa by $1$. In
Chapter~\ref{ch:XII} the excess emptied the boundary quotient; here it displaces the equilibrium
range out of reach of the obstruction. It is one fact with two consequences.
\end{remark}

\section{Sparse sets for $\Delta$}

\begin{theorem}[Anomalous sets exist]\label{XXI:thm:sparse}
Let $\Delta\in S_{12}(SL_2(\Z))$ with $\tau(p)=2p^{11/2}\cos\theta_p$. Assuming the Sato--Tate
equidistribution for $\Delta$ --- a theorem of Barnet-Lamb, Geraghty, Harris and Taylor \cite{BLGHT} --- for
every $\delta>0$ the set $\{p:\theta_p\le\delta\}$ has positive density. Choosing finite
blocks $S_j\subseteq\{p:\theta_p\le2^{-j}\}$ with $1\le\sum_{p\in S_j}p^{-1}\le2$ and putting
$S_\Delta=\bigsqcup_jS_j$ gives
\[
\sum_{p\in S_\Delta}p^{-1}=\infty,\qquad \Xi_1(\pi_\Delta,S_\Delta)\le\sum_j2\cdot4^{-j}\cdot2<\infty .
\]
\end{theorem}

\begin{proof}
$8\sin^2(\theta/2)\le2\theta^2$, so the $j$-th block contributes at most
$2\cdot4^{-j}\sum_{p\in S_j}p^{-1}\le4\cdot4^{-j}$. This is Theorem 5.3 of Chapter~\ref{ch:VI} with
$\Delta$ in place of an elliptic curve.
\end{proof}

\begin{remark}[But not symmetry breaking]\label{XXI:rem:notbreaking}
Theorem~\ref{XXI:thm:sparse} produces a set on which the analytic obstruction is finite. It does
not produce a broken phase, for two reasons already recorded. By
Remark 5.4 of Chapter~\ref{ch:VI} an obstruction group is not a broken phase until an algebra realizes
it; and by Theorem~\ref{XXI:thm:gap} the temperature $\beta=1$ at which the construction operates
is one at which this $GL_2$ system has no equilibrium states, $\beta_c(S_\Delta)=1$ forcing
Gibbs states only above $2$. The residual symmetry subgroup $H_\Delta\subseteq SU(2)$ asked
for is therefore not defined: there is no simplex on which it could act.
\end{remark}

\section{Assessment}

\[
\begin{array}{lll}
\text{Proposal} & \text{Status} & \text{Reference}\\\hline
C(Y)\rtimes M_2(\Z_S)^+ & \textbf{no: }|P^\times|=\infty & \text{Prop.~1.1}\\
\|I_2-U_p\|^2_\HS=8\sin^2(\theta_p/2) & \text{true} & \text{Lem.~2.1}\\
\Xi_\beta\ \text{involves}\ \Sym^2 & \textbf{false: spurious term} & \text{Rem.~2.3}\\
\Xi_\beta=4\log\zeta_S-2\Re\log L_S+O(1) & \textbf{true} & \text{Thm.~2.2}\\
\Xi_1(\pi_f,P_\Q)=\infty & \text{true, Jacquet--Shalika} & \text{Cor.~2.4}\\
\text{uniqueness on }(0,1] & \textbf{vacuous: no KMS states there} & \text{Cor.~3.3}\\
\text{obstruction region} & (0,\beta_c(S)],\ \beta_c\le1 & \text{Thm.~3.2}\\
\text{equilibrium region} & (1+\beta_c(S),\infty) & \text{Thm.~3.2}\\
\text{the two overlap} & \textbf{never; gap exactly }1 & \text{Thm.~3.2}\\
\text{sparse }S_\Delta\ \text{with}\ \Xi_1<\infty & \text{true, under Sato--Tate} & \text{Thm.~4.1}\\
\text{residual symmetry}\ H_\Delta\subseteq SU(2) & \textbf{not defined} & \text{Rem.~4.2}
\end{array}
\]

The arithmetic side of the proposal is sound and largely already proved in
Chapter~\ref{ch:VI}: the Hilbert--Schmidt series, its $L$-function form once the spurious
symmetric square is removed, the divergence at $\beta=1$ by Jacquet--Shalika, and the
construction of Sato--Tate anomalous sets. What fails is the junction with statistical
mechanics, and it fails quantitatively rather than vaguely. The obstruction is a Dirichlet
series in $p^{-\beta}$, with abscissa at most $1$; the $GL_2$ partition function counts
sublattices, whose local factor $(1-p^{1-\beta})^{-1}$ moves the abscissa to
$1+\beta_c(S)$. The gap is exactly $1$ and it is the same off-by-one of $\mathbb P^1$ that
Chapter~\ref{ch:XII} showed to be universal.

Two ways out suggest themselves and neither is free. One could weight the dynamics
differently, replacing $|\det g|$ by another multiplicative norm; but by
Proposition 1.4 of Chapter~\ref{ch:X} norms factor through the abelianization, and on $GL_2$ orbits the
abelianization is generated by $\det$, so any admissible norm is $|\det|^t$ and merely
rescales $\beta$, leaving the ratio of the two abscissae fixed at
$(1+\beta_c)/\beta_c$ rather than closing the gap. Or one could keep the $GL_1$
normalization and attach the automorphic data as a coefficient system rather than as the
acting semigroup --- which is the route of \S6 of Chapter~\ref{ch:VI}, where the Kakutani analysis
applies verbatim and only the existence of a tractable expectation was left open. We would
pursue the second.

